\subsection{Neutrino data analysis of extra-dimensional theories with massive bulk fields --- arXiv:2508.04274}

\textbf{Authors:} Philipp Eller, Manuel Ettengruber, Alan Zander

\paragraph{主な主張}
MINOS/MINOS+・KamLAND・Daya Bayの結合解析でbulk Dirac質量を持つLEDの兆候はなく、半径制約がbulk質量とYukawa結合の仮定に強く依存することを示す\cite{Eller:2025MassiveBulk}。両者を同時に自由化した解析では、探索範囲内で排除領域が得られなかった。

\paragraph{新規性}
質量を持つbulk fermionのLED+とDark Dimensionを共通の尤度解析で扱い、基本重力スケールと半径の関係の違いも含めて制約の模型依存性を定量化した。

\paragraph{位置付け}
最小LEDの半径上限を一般の余剰次元模型に直ちに適用できないことを示す近年の解析であり、幾何学の効果と局在・結合の効果を分離する上で重要である。
